\chapter{Background}
\label{ch:background}

\section{APR Pipeline}
\label{sec:bg-pipeline}

Figure~\ref{fig:bg-pipeline} places validation in the full repair pipeline.
Fault localisation finds the code most likely to hold the bug. The generator
proposes a candidate patch there, and the validator builds it and runs the
test suite. A plausible candidate goes on to a human or automatic correctness
check; a failing one, in an LLM-based loop, returns its failure to the model,
which proposes the next candidate.

\begin{figure}[htbp]
\centering
\resizebox{0.98\textwidth}{!}{%
\begin{tikzpicture}[
  font=\small,
  st/.style={draw=black!55,rounded corners=2pt,align=center,inner sep=4pt,
             minimum height=11mm,text width=21mm,fill=black!3},
  vs/.style={draw=bkblue,line width=0.9pt,rounded corners=2pt,align=center,
             inner sep=4pt,minimum height=11mm,text width=21mm,fill=bkblue!10},
  arr/.style={-{Stealth[length=2.4mm]},draw=black!60}
]
\node (fl)  at (0,0)    [st] {fault\\ localisation};
\node (gen) at (3.2,0)  [st] {patch\\ generation};
\node (bld) at (6.4,0)  [vs] {build};
\node (run) at (9.6,0)  [vs] {run tests};
\node (pl)  at (13.8,0) [st] {plausible\\ patch};
\node (ok)  at (17.0,0) [st] {correctness\\ check};
\draw[arr] (fl) -- (gen);
\draw[arr] (gen) -- (bld);
\draw[arr] (bld) -- (run);
\draw[arr] (run) -- node[above,pos=0.6,font=\scriptsize]{all pass} (pl);
\draw[arr] (pl) -- (ok);
\draw[arr] (run.south) -- ++(0,-0.75) -| node[pos=0.25,below,font=\scriptsize]
  {a test fails: feedback for the next candidate} (gen.south);
\node[draw=bkblue,dashed,rounded corners=3pt,fit=(bld)(run),inner sep=6pt,
      label={[font=\scriptsize\itshape,text=bkblue]above:validation}] {};
\end{tikzpicture}}
\caption[The generate-and-validate pipeline]{The generate-and-validate
pipeline. Validation (highlighted) runs once per candidate.}
\label{fig:bg-pipeline}
\end{figure}

Part of each validation is overhead (recompiling and starting the runtime),
which on the JVM can be removed by validating in one long-lived
JVM~\citep{chen2021uniapr,xiao2024expressapr}; the rest is running tests.

\section{Cost Model of a Repair Session}
\label{sec:bg-cost}

This section makes precise the first question of the thesis: where the cost
of a repair session goes. A session runs rounds $i = 1, \dots, R$; in each round the model
proposes one candidate $p_i$ and the validator judges it. The session ends at
the first plausible candidate or when the round budget $R_{\max}$ is spent. Its
cost is
\begin{equation}
E \;=\; \sum_{i=1}^{R} G_i \;+\; \sum_{i=1}^{R} V_i ,
\label{eq:bg-e2e}
\end{equation}
where $G_i$ is the cost of the model call of round $i$ and $V_i$ everything it
costs to validate its candidate: applying the patch, building it, starting the
runtime and running the tests. Cost is measured in three units, each reported
on both sides: wall-clock time, which is what a user waits for; compute (CPU
seconds for validation, tokens for generation); and money, from a fixed price
list. The model has two layers: what the session spent, and what each family of
techniques could have saved.

\paragraph{What the session spent.} A candidate is a \emph{repeat} if its
normalised text equals that of a candidate already proposed in the same
session. Splitting the rounds by whether their candidate is a repeat, and the
validated candidates by their verdict, gives five components, each read
directly from the session:
\begin{equation}
E \;=\; \underbrace{G_{\mathrm{new}} + G_{\mathrm{rep}}}_{\text{generation}}
\;+\; \underbrace{V_{\mathrm{rep}} + V_{\mathrm{rej}} + V_{\mathrm{acc}}}_{\text{validation}} ,
\label{eq:bg-decomp}
\end{equation}
where $G_{\mathrm{new}}$ and $G_{\mathrm{rep}}$ are the generation of new and of
repeated candidates; $V_{\mathrm{rep}}$ is the validation of repeats, which a
verdict cache removes; $V_{\mathrm{rej}}$ the validation of new candidates that
were rejected; and $V_{\mathrm{acc}}$ the validation of the accepted candidate,
which runs the whole suite in any order.

\paragraph{What each family could save.} Each family of techniques
(Section~\ref{sec:rel-taxonomy}) acts on part of Equation~\eqref{eq:bg-decomp}, and
that part bounds what it can save (Table~\ref{tab:bg-headroom}). Two bounds need
a definition. The \emph{ordering headroom} is $V_{\mathrm{rej}} -
V_{\mathrm{rej}}^{*}$, where $V_{\mathrm{rej}}^{*}$ is what rejecting the same
candidates would have cost had each been rejected by its cheapest rejecting test, which is the least any
order can pay (Section~\ref{sec:bg-problem}). The
\emph{start-up share} is the part of validation spent building candidates and
starting the runtime rather than executing tests. Unlike the five components,
the bounds are counterfactual; all of them are computed after the session from
the outcome and cost of every test on every distinct candidate, without running
any technique.

\begin{table}[htbp]
\centering
\caption[What each family of techniques can save]{What each family of
techniques can save at most, in the terms of Equation~\eqref{eq:bg-decomp}.}
\label{tab:bg-headroom}
\footnotesize
\begin{tabularx}{\textwidth}{@{}>{\raggedright\arraybackslash}p{0.32\textwidth} >{\raggedright\arraybackslash}X@{}}
\toprule
Family & Upper bound on its saving\\
\midrule
Verdict cache (patch reduction) & $V_{\mathrm{rep}}$\\
Test prioritisation & the ordering headroom, $V_{\mathrm{rej}} - V_{\mathrm{rej}}^{*}$\\
Test selection & the cost of tests that do not reach the change; zero where every test runs the whole program\\
Execution acceleration & the start-up share of validation\\
Generation-side reduction & $G_{\mathrm{rep}}$, plus the validation of the rounds saved\\
\bottomrule
\end{tabularx}
\end{table}

The share of each component in $E$ is not a property of a technique but of a
setting. The generation side scales with the latency and price of the model;
the start-up share with the runtime and the build; the rest of validation with
the number of tests, their cost and its spread. Two consequences follow. A
technique can save a large fraction of its own component and a small fraction
of $E$; and the same technique can be worth deploying in one setting and not in
another.

\section{Test Ordering under Early Stopping}
\label{sec:bg-problem}

This section states the problem test ordering solves: it concerns the test-execution part of $V_i$ in
Equation~\eqref{eq:bg-e2e}. A session produces candidates $p_1, p_2, \dots$ in order and has a test suite
$T$. For a candidate $p$ and test $t$, let $o(p,t) = 1$ if $t$ rejects $p$ and
$0$ otherwise, and let $c(p,t) > 0$ be the cost of running it (CPU time unless stated
otherwise). Neither is
known before the test runs. A candidate is plausible when no test rejects it.

\begin{definition}[Validation with early stopping]
\label{def:strategy}
For each candidate $p$, the validator runs the tests in an order $\sigma_p$ and
stops at the first test that rejects $p$. The order for a candidate may depend
on everything observed on earlier candidates.
\end{definition}

The validation cost of a session is
\begin{equation}
C \;=\; \sum_{p} \; \sum_{j=1}^{r_p} c\bigl(p, \sigma_p(j)\bigr),
\label{eq:bg-cost}
\end{equation}
where $r_p$ is the position of the first rejecting test, or $|T|$ if none
rejects, and the sum runs over the candidates actually validated: a repeat
whose verdict is served from a cache contributes nothing. In the terms of
Equation~\eqref{eq:bg-decomp}, $C$ is the test-execution part of
$V_{\mathrm{rej}} + V_{\mathrm{acc}}$ with a verdict cache, and of
$V_{\mathrm{rep}}$ as well without one. A technique is \emph{safe} if it declares a candidate plausible
exactly when it is. Reordering is safe whenever tests are deterministic and
their outcomes do not depend on the order in which they run, since every test
still runs before a candidate is accepted.

\paragraph{Lower bound.} If all outcomes were known,
the cheapest cost would be
\begin{equation}
C^{*} \;=\;
\sum_{p \text{ invalid}} \; \min_{t \,:\, o(p,t)=1} c(p,t)
\;+\;
\sum_{p \text{ plausible}} \; \sum_{t \in T} c(p,t).
\label{eq:bg-hindsight}
\end{equation}
The second term is fixed: a plausible candidate runs the whole suite in any order. Ordering acts only on the first term, the cost of rejecting invalid
candidates. In a session that stops at the first plausible candidate, every
candidate but the last is invalid by construction. Search-based tools generate
many invalid candidates, so there the first term dominates. An LLM-based loop validates few distinct candidates per session,
so the second term (one full suite) can be of the same order as the first, and the room left for ordering is correspondingly smaller. Which term dominates is measured in the experiments (Chapter~\ref{ch:experiments}).

\paragraph{Optimal order.} The validator does not know the outcomes in
advance, so it has to estimate them. If each test rejected independently with
probability $q_t$ at cost $c_t$, the best order is a classical result: the
least-cost testing sequence~\citep{mitten1960}, also derived for satisficing
search~\citep{simon1975satisficing}. It is restated here, with
the exchange argument, because the rest of the thesis estimates its two
quantities.

\begin{proposition}[Rate-over-cost order~\citep{mitten1960,simon1975satisficing}]
\label{prop:ratio}
Under independent rejection probabilities $q_t$ and costs $c_t$, running tests
in decreasing $q_t / c_t$ minimises the expected cost of rejecting an invalid
candidate.
\end{proposition}
\begin{proof}
For adjacent tests $a$ and $b$, running $a$ first costs $c_a + (1-q_a)\,c_b$
in expectation and running $b$ first costs $c_b + (1-q_b)\,c_a$. So $a$ should
go first exactly when $q_a/c_a > q_b/c_b$. Any other order contains an adjacent pair whose swap lowers the expected cost.
\end{proof}

\paragraph{Example.} Take two test classes. Class $A$ rejects a wrong patch
with probability $0.5$ and takes $10$ seconds; class $B$ rejects it with
probability $0.2$ and takes $1$ second. Running $A$ first costs on average
$10 + 0.5 \times 1 = 10.5$ seconds, since $B$ runs only when $A$ passes;
running $B$ first costs $1 + 0.8 \times 10 = 9$ seconds. The rule gives the same
answer: $q/c$ is $0.05$ for $A$ and $0.2$ for $B$, so $B$ goes first even though
$A$ fails more often. A failure count would put $A$ first.

A failure count, as in \frtp, estimates neither $q_t$ nor $c_t$.
Proposition~\ref{prop:ratio} is the target; the contribution of this thesis is
not the rule but estimating its two quantities online, from what validation
reveals.

\paragraph{Target rate.} When rejections are correlated, what matters for a test
at a given position is its rejection probability \emph{given} that every test
before it has passed. A record built under early stopping estimates a rate of this kind (rejections among the candidates that reached
the test), but averaged over whatever tests happened to precede it, so it equals neither the
marginal rate nor the conditional rate of any one order. The thesis takes the
rate conditional on reaching the test as the target, and will measure how far
the record departs from it (Chapter~\ref{ch:experiments}).

\section{Regression Test Prioritisation}
\label{sec:bg-regression}

Ordering tests is an old problem in regression testing, where a changed program
is re-tested against an existing suite.
\citet{rothermel2001prioritizing} introduced the two greedy strategies that
later studies, including the study of validation in repair by
\citet{lou2024aprregression}, use as baselines: \emph{total} (\tptot) ranks tests by how much code they
cover, and \emph{additional} (\tpadd) by how much they cover that earlier tests
did not. An order is scored by how early it detects faults:
\begin{equation}
\mathrm{APFD} \;=\; 1 - \frac{\mathit{TF}_1 + \dots + \mathit{TF}_m}{n\,m} + \frac{1}{2n},
\label{eq:apfd}
\end{equation}
where $n$ is the number of tests, $m$ the number of faults and $\mathit{TF}_i$
the position of the first test revealing fault $i$.
\citet{elbaum2001cost} extended it to APFD$_C$, which weighs each test by its
cost.

Two differences separate patch validation from this setting. Under early
stopping only the \emph{first} failure of each candidate matters, so the relevant score is the cost paid before it (Equation~\eqref{eq:bg-cost})
rather than how early all faults are found. And in a repair loop the
``program'' changes with every candidate, so an order can learn from one
candidate to the next.

\section{Censored Test Outcomes}
\label{sec:bg-censoring}

Picture the outcomes of a session as a matrix: one row per candidate, one
column per test. An order learns from this matrix, but under early stopping it
sees only part of it. For each invalid candidate it observes a few passes and
one rejection, and nothing about the tests after the rejection. The record is
\emph{censored}, and by construction not at random: a test is observed only
when every test ranked before it has passed, so the tests ranked first are
observed most often. Four consequences follow.

\begin{itemize}[leftmargin=*,itemsep=3pt]
\item \textbf{A count is not a rate.} A count of rejections gives zero both to
a test never run and to a test run many times without rejecting. A rate
(rejections over runs) separates them, but is estimated from few observations
for the tests ranked low.
\item \textbf{Joint rejections are never observed.} Validation stops at the
first rejection, so the record of a candidate never shows two tests rejecting
it. Whether two tests reject the same candidates, a relation that selection and
redundancy arguments rely on, cannot be read from the record,
and a rule that needs it cannot be fed by the session's own runs.
\item \textbf{Cost estimates need repeated runs.} One rejection already shows
that a test can reject; a cost estimate is an average of timings and becomes
reliable only after a test has run several times. Each test's cost on the
faulty program is measured once before repair, but its cost on a candidate is
observed only when it runs on that candidate.
\item \textbf{The rates need not be stationary.} In a conversational loop the
failing test is shown to the model, which then aims at it; a test that has just
been reported may reject the next candidate less often. A rate pooled over the
whole session assumes otherwise.
\end{itemize}

A history-based order is therefore learned from data that its own early
stopping withholds.

\paragraph{The cascade view.} This is the observation model of \emph{cascading
bandits}~\citep{kveton2015cascading}, a learning-to-rank setting in which a user
scans a ranked list and stops at the first attractive item; here the validator
scans the ordered tests and stops at the first rejection. Two facts carry over.
First, if outcomes are independent, whether a test is reached does not depend
on its own outcome, so its rate among the runs that reached it is unbiased;
with correlated outcomes the same rate is the rate conditional on reaching,
the target of Section~\ref{sec:bg-problem}. Censoring thus costs observations
more than accuracy. Second, that cost falls on the tests ranked low: a test
whose rate is underestimated is rarely reached, so its estimate is rarely
corrected. Reweighting the record by the inverse probability of reaching each
test does not help, because a deterministic order reaches a test with
probability zero or one. An order obtains the missing observations only by
\emph{exploring}: running an uncertain test earlier than its estimate alone
would.

\section{Benchmark}
\label{sec:bg-benchmarks}

Repair tools are evaluated on benchmarks of real bugs, each shipped with a
test suite that exposes the bug. This thesis uses \textbf{Defects4J}, the
common benchmark of Java repair~\citep{just2014defects4j}: real bugs from
open-source projects, each with the developer's fix and at least one test that
fails before the fix and passes after it. UniAPR~\citep{chen2021uniapr} and the
study of test selection in repair~\citep{lou2024aprregression} were also
evaluated on it. Other benchmarks differ in what a test is: in
ConDefects~\citep{wu2024condefects} a test is one run of a whole contest
program, and in SWE-bench~\citep{jimenez2024swebench} a test of a large Python
repository. Two properties of a benchmark decide how much a validation
technique can help: how fine-grained a test is, which decides whether test
selection has anything to remove, and how much test costs differ within one
bug, which decides whether the order of tests matters in seconds.
